\section{Results IV: Interaction-Specific Assumptions}

Several experiments directly test assumptions that are plausible at the level of product design
but are not guaranteed by recognizer accuracy alone.

\subsection{Beam width}

For the default \texttt{base.en} checkpoint, greedy decoding has a worse point estimate than
beam 5 on both splits. On the clean split, beam 1 produces 4.39\% WER compared with 4.01\%;
the paired difference is 0.386 WER points with a 95\% interval of 0.069 to 0.726 and an
unadjusted pairwise $p=0.0192$. On the hard split, beam 1 produces 10.56\% compared with
9.46\%; the paired difference is 1.102 points with a 95\% interval of 0.305 to 2.196 and
$p=0.001$. A conservative Bonferroni correction over the five archived beam comparisons within
each split gives $\alpha=0.01$: under that reading the hard-split beam-1 versus beam-5 result
survives, while the clean-split comparison does not.

The evidence does not support the stronger rule that more beam is always better. On the hard
split, beams 2, 3, 5, and 8 are not shown to differ meaningfully from one another in the paired
comparisons against beam 2. On \texttt{small.en} clean speech, the beam-1 point estimate is
slightly better than beam 5, but the paired interval crosses zero. The practical distinction in
the measured default-model condition is therefore greedy decoding versus a modest beam, not an
unbounded preference for larger beam widths.

\subsection{Synthetic pre-speech padding}

The onset experiment clips 0, 40, 120, or 240 ms of real speech and prepends 0, 100, 300, 600,
or 1000 ms of synthetic silence. The primary outcome is first-word correctness on 200
\texttt{test-clean} utterances.

No tested lead-in comparison survives multiplicity correction as an improvement over no
synthetic lead-in. When the onset is intact, the shipped 300 ms padding does not establish a
measurable first-word benefit. Once real speech has been clipped, prepending silence cannot
reconstruct the missing acoustic evidence. Two replicated uncorrected cells at 120 ms clipping
move in the harmful direction for very long lead-ins.

This negative result clarifies the privacy/interaction trade-off. Recovering speech that occurred
before explicit activation would require retaining real pre-activation microphone audio. Synthetic
silence preserves the stronger hold-to-talk privacy boundary, but cannot provide that recovery.

\subsection{CPU streaming}

Streaming partial decoding is useful only when the rolling recognizer can stay sufficiently far
ahead of incoming audio. We report the committed manuscript-grade artifact
\texttt{paper/results/streaming.json}: 15 real-time-fed utterances, 100 ms audio chunks, and a
300 ms partial-decode interval. Unlike the engine matrix, this artifact was produced on the
reference laptop (20-thread Intel Core i7-1370P) rather than the Azure host, so its absolute
latencies belong to that machine; the batch-versus-streaming contrast within it is paired on the
same machine and the same utterances. For \texttt{tiny.en}, median final speech-end-to-text
latency increases from 1.377 s in batch mode to 2.371 s with streaming, while a median 59.33\%
of final text is visible at release.

For \texttt{base.en}, median final speech-end-to-text latency increases from 4.024 s to 7.612 s.
Median visible text at release is 0\%, and 12 of 15 utterances produce no confirmed partial
before release. The partial loop consumes CPU but often fails to deliver a useful preview before
the final decode is required.

The appropriate conclusion is therefore conditional. Streaming can reduce perceived latency for
a sufficiently fast checkpoint, but it is not itself a general latency optimization on CPU.